\subsection{A reverse mirror transition for \texorpdfstring{$X_9$}{X9}}
\label{ssec:x9reverse}

For $X_9$ the projective resolutions of Theorem~\ref{thm:x9global}
admit an explicit comparison with the Batyrev mirror family of a
smoothing. We first identify the smoothing of the particular nodal
partial resolution $X'_9$ used in this paper.

\begin{theorem}\label{thm:x9reverse}
The threefold $X'_9$ of Section~\ref{ssec:x9} has a proper flat analytic
smoothing whose noncentral fibres are smooth projective
anticanonical hypersurfaces from the toric smoothing family of
\cite{Paper5}. For general $c\in\mathcal D\cap S_4$, the projective
small resolution of $Y_f$ in Theorem~\ref{thm:x9global} is birational to
a smooth projective member of the Batyrev mirror family of these
smoothing fibres.
\end{theorem}

Thus the reverse transition is established here in the sense of
Batyrev families, allowing birational Calabi--Yau models for the resolved
mirror. The theorem concerns a proper analytic smoothing with projective
fibres; it does not require a relatively projective total space.

\begin{proof}[Proof of Theorem~\ref{thm:x9reverse}]
\emph{A simultaneous modification of the degree-seven smoothing.}

The degree-seven cone has the height-one polygon
\[
 \Pi=\operatorname{conv}\{(0,0),(1,0),(2,1),(1,2),(0,1)\}
       =[0,(1,1)]+\operatorname{conv}\{0,(1,0),(0,1)\}.
\]
In coordinates $(p,q,a,b)$ on $\Z^4$, its Cayley cone has rays
\begin{equation}\label{eq:x9cayleyrays}
\begin{gathered}
 A_0=(0,0,1,0),\quad A_1=(1,1,1,0),\\
 B_0=(0,0,0,1),\quad B_1=(1,0,0,1),\quad B_2=(0,1,0,1),
\end{gathered}
\end{equation}
with relation $-A_0+A_1+2B_0-B_1-B_2=0$. Let $\mathcal C$ be the
affine toric fourfold of this cone. Altmann's component construction
\cite[\S8.1]{Altmann97} gives the reduced smoothing family as the map
$\mathcal C\to\C_s$ defined by the difference of its two height
characters. The full local base is
$\C\{s,w\}/(w^2,sw)$ \cite[\S9.1(v)]{Altmann97}; in particular every
map from a reduced disc factors through this reduced line.

Subdivide the cone into
\begin{equation}\label{eq:x9cayleycones}
 \langle A_1,B_0,B_1,B_2\rangle,\quad
 \langle A_0,A_1,B_0,B_2\rangle,\quad
 \langle A_0,A_1,B_0,B_1\rangle.
\end{equation}
The determinants are $1,-1,1$. Heights $(1,0,0,1,1)$ on the five
rays give this regular subdivision: on each cone the omitted ray lies
three units above its linear height. The resulting toric morphism is
\[
 \rho:\mathcal V=\operatorname{Tot}_{\mathbb P^2}
                \bigl(\mathcal O(-1)\oplus\mathcal O(-2)\bigr)
                 \longrightarrow\mathcal C.
\]
Indeed the relation in \eqref{eq:x9cayleyrays} gives quotient weights
$(-1,1,2,-1,-1)$ on $(u_0,u_1,v_0,v_1,v_2)$, with stable set
$(u_0,v_1,v_2)\ne0$. Write $[X_0:X_1:X_2]=[u_0:v_1:v_2]$ and
fibre coordinates $U=u_1$, $V=v_0$. The equivalence is
\[
 (X,U,V)\sim(\lambda X,\lambda^{-1}U,\lambda^{-2}V).
\]
The invariant functions are generated by $X_iU$ and $X_iX_jV$.
The tautological bundle inclusions give a closed immersion
$\mathcal V\hookrightarrow\mathbb P^2\times\C^9$, so the map
to its affine invariant image is projective. Its fibre over the vertex
is exactly the zero section. Off that vertex it is an isomorphism:
if $X_iU\ne0$, the ratios $(X_jU)/(X_iU)$ recover the base; if all
these vanish, some $X_i^2V$ is nonzero and the ratios
$(X_iX_jV)/(X_i^2V)$ recover it. The invariant relations ensure that
these opens cover the complement of the vertex, and the same formulas
recover the fibre coordinates regularly.

The height characters are $u_0u_1$ and $v_0v_1v_2$. Hence the parameter
on $\mathcal V$ is
\begin{equation}\label{eq:x9simultaneous}
 s=X_0U-X_1X_2V.
\end{equation}
It is flat over $\C_s$: on affine charts, the nonconstant function $s$
makes the coordinate ring torsion-free over the principal ideal domain
$\C[s]$. On $X_1\ne0$, with
$a=X_0/X_1$, $b=X_2/X_1$, $u=X_1U$, $v=X_1^2V$, the equation is
$s=au-bv$. Its unique critical point is the chart origin. The
$X_2\ne0$ chart gives the other critical point, and on $X_0\ne0$ the
equation is $s=u-pqv$, with no critical point. Thus the central fibre
has exactly two nodes, at $[0:1:0]$ and $[0:0:1]$ on the zero section,
and all noncentral fibres are smooth. The contraction is an isomorphism
on every noncentral fibre.

The central fibre is the precise partial resolution used in $X'_9$.
To see this, intersect the cones in \eqref{eq:x9cayleycones} with
$a=b=1$. The three cells are
\[
\begin{gathered}
 \{(1,1),(2,1),(1,2)\},\\
 \{(0,0),(1,1),(1,2),(0,1)\},\qquad
 \{(0,0),(1,0),(2,1),(1,1)\}.
\end{gathered}
\]
The affine lattice map $(p,q)\mapsto(1-p,1-q)$ takes them to
$T_3,P_4,P_1$ of Lemma~\ref{lem:subdiv}, and takes the new central
ray to the pentagon centre. The central charts are one affine
three-space and two normal node charts. Thus the binomial central
fibre is already normal and is the toric variety with this intersection
fan. Its exceptional divisor is $\mathbb P^2$, with precisely the two
nodes prescribed by our pentagon subdivision. The total space has
trivial canonical bundle, since
$\det(\mathcal O(-1)\oplus\mathcal O(-2))=K_{\mathbb P^2}$, and
the central map is crepant.

\emph{The original-side smoothing fibres.}

The ambient family in \cite[Theorem~11.8]{Paper5} gives a smoothing
of the chosen mixed threefold $X_9$, which has two nodes and one
degree-seven cone point. To keep this precise central hypersurface,
let $g_0$ be its anticanonical section. By
\cite[Proposition~11.7]{Paper5}, the relative section bundle is locally
free of rank $162$ and commutes with base change. Trivialize it on a
disc with coordinate $t$, and extend $g_0$ by
$g(t)=g_0+t h$. For a fixed $t_1\ne0$, varying $h$ gives every section
on the fibre at $t_1$. Smooth sections form a dense open subset, so
choose $h$ giving a smooth hypersurface there and not proportional to
$g_0$. Then $g(t)$ is nonzero throughout the original disc. The
fibres of the ambient family are integral, so the fibrewise
regular-section criterion makes $\{g(t)=0\}$ flat over the disc.
It is proper. The image of its nonsmooth locus is a proper analytic
subset of the disc, since the fibre at $t_1$ is smooth. After a further
shrinking, all noncentral fibres are therefore smooth, while the
central fibre remains the chosen $X_9$.

Choose the central analytic identification of the cone point compatibly with the
partial resolution $X'_9\to X_9$. In the toric hypersurface charts,
this is the identification obtained by eliminating the transverse
torus coordinate. Analytic semiuniversality identifies the actual
local family with the pullback of the degree-seven versal family along
a classifying arc. The disc is reduced, so its $w$-coordinate vanishes;
the arc therefore factors through the reduced family just described.
Its parameter $s(t)$ is nonzero for $t\ne0$ after shrinking, because
the arc is a smoothing. No assumption on $s'(0)$ is needed.

Pull back $\rho$ along this arc and restrict it over a representative
of the identified neighbourhood in the total family. Flatness survives
base change. Its only exceptional locus maps to the central cone point,
so it glues to the identity on the complement of that point. The glued
space is Hausdorff: points with distinct target images are separated
downstairs, while points with the same target image are separated in
one inverse image of the target open cover. Properness is local on the
target, and flatness is local on the source. We obtain a proper flat
analytic family with central fibre $X'_9$ and noncentral fibres
isomorphic to those of the chosen anticanonical family. The two original nodes are
unchanged in the central fibre and are smoothed by the same family.

We recall the toric ambient of those smooth fibres, to fix the mirror
comparison. Its fan polytope $\nabla$ has vertices
\begin{equation}\label{eq:x9smoothrays}
\begin{aligned}
 v_0&=(-1,-3,-1,0),&v_1&=(-1,-3,0,-1),\\
 v_2&=(-1,-1,0,0),&v_3&=(-1,0,0,0),\\
 v_4&=(0,1,0,0),&v_5&=(1,-1,0,0),&v_6&=(1,3,1,1).
\end{aligned}
\end{equation}
The maximal cones, as sets of indices, are
\begin{equation}\label{eq:x9smoothcones}
 1456,\ 1346,\ 0456,\ 0346,\ 0145,\ 0134,\ 01256,\ 1236,\ 0236,\ 0123.
\end{equation}
This is the Gorenstein Fano fourfold computed in \cite{Paper5}. Its
only singular point is the fixed point of $01256$, which a general
anticanonical member avoids. The relative anticanonical sections form a
rank-$162$ vector bundle, commute with base change and generate the
relative line bundle. Consequently the noncentral hypersurfaces occur
in the complete anticanonical family, whose Newton polytope is
$\nabla^\circ$. Their Batyrev partner has Newton polytope $\nabla$.

\emph{The birational mirror family.}

On the dense torus, use the parameters of
\eqref{eq:x9factorparameters} and the unimodular coordinates
\[
 a=x,\qquad b=z/x^2,\qquad c=w/x^2,\qquad t=y/x.
\]
Put
\begin{equation}\label{eq:x9mutationfactors}
\begin{split}
 g&=1+ra^4bc,\\
 h&=A/a+B/a^2+C/(a^4b)+D_0/(a^4c),\\
 k&=c_0a+c_{11}+c_9/a.
\end{split}
\end{equation}
Equation~\eqref{eq:x9factorF}, divided by $x^2t$, is
$c_1t+k+gh/t=0$. The rational substitution $t=g\zeta$ gives
\begin{equation}\label{eq:x9mutant}
 c_1g\zeta+k+h/\zeta=0.
\end{equation}
It has inverse $\zeta=t/g$ on $g\ne0$, and hence identifies dense
opens of the two hypersurfaces. The nine possible exponents of
\eqref{eq:x9mutant}, in $(a,b,c,\zeta)$, are
\begin{equation}\label{eq:x9mutantsupport}
\begin{gathered}
 (0,0,0,1),\ (4,1,1,1),\ (1,0,0,0),\ (-1,0,0,0),\ (0,0,0,0),\\
 (-1,0,0,-1),\ (-2,0,0,-1),\ (-4,-1,0,-1),\ (-4,0,-1,-1).
\end{gathered}
\end{equation}
Their convex hull has seven vertices and exactly these nine lattice
points. The lattice isomorphism
\[
 (u_1,u_2,u_3,u_4)\longmapsto(u_4,u_1-u_4,u_2,u_3)
\]
has determinant $-1$ and sends its vertices to
\eqref{eq:x9smoothrays}. The ten facets have exactly the index sets
in \eqref{eq:x9smoothcones}. This identifies its Newton polytope with
$\nabla$, including the lattice, not only the face combinatorics.

The coefficient map is equally explicit. The nine parameters on the
open part of $\mathcal D$ map to the nine coefficients of
\eqref{eq:x9mutant} as
\begin{equation}\label{eq:x9coefficientmap}
 (c_1,r,c_0,c_{11},c_9,A,B,C,D_0)\longmapsto
 (c_1,c_1r,c_0,c_{11},c_9,A,B,C,D_0),
\end{equation}
up to the displayed ordering of monomials. This is an isomorphism of
coefficient opens, with inverse $r=\text{second}/\text{first}$.
In particular the constant coefficient $c_{11}$ is unrestricted.
The compact four-node good locus and the inverse image of the
Newton-nondegenerate locus are nonempty Zariski-open subsets of the
irreducible $\mathcal D$, so intersect. For a general member of the latter, Batyrev's maximal projective
crepant partial desingularization is a smooth projective threefold
by \cite[Thm.~4.2.2, Cor.~4.2.3]{Batyrev}.
Its function field agrees with that of the projective small resolution
of $Y_f$ from Lemma~\ref{lem:x9surface}. This proves
Theorem~\ref{thm:x9reverse}.
\end{proof}

The torus substitution also preserves the logarithmic volume form,
since $d\log t=d\log\zeta+d\log g$ and the latter summand involves
only $a,b,c$. This identity on the common open does not assert a
comparison of compact quantum invariants. At special points of
$\mathcal D\cap S_4$, the Laurent polynomial
\eqref{eq:x9mutant} may fail to be Newton nondegenerate; the theorem's
mirror identification is stated for general members.
